\section{Emergent Ontology}
\label{sec:ontology}

Human cognition forms categories organically through experience---the
Yogac\=ara concept of \emph{vikalpa} (conceptual construction) describes how
the mind labels and structures percepts without explicit instruction.
\alaya{}'s emergent ontology implements a pragmatic approximation: categories
emerge from clustering patterns in the semantic store during lifecycle
operations, providing a lightweight taxonomic structure that requires no
manual curation.

\subsection{Design Principles}

The emergent ontology follows three principles:

\begin{enumerate}[nosep]
    \item \textbf{Dual-signal evidence:} Categories are supported by both
    embedding similarity (content-level) and graph topology (usage-level).
    Neither signal alone is sufficient---embeddings capture semantic
    similarity while graph links capture co-access patterns that reflect
    the user's mental model.

    \item \textbf{Separation of assignment and discovery:} Consolidation
    assigns new nodes to \emph{existing} categories (zero LLM cost);
    transformation discovers \emph{new} categories from uncategorized nodes
    (one cheap LLM call per category for labeling, optional).

    \item \textbf{Graceful degradation:} Without embeddings, the system
    falls back to graph-only clustering (shared Hebbian link overlap).
    Without an LLM, categories receive placeholder labels derived from
    prototype node content. The system never fails due to missing optional
    capabilities.
\end{enumerate}

\subsection{Data Model}

Each category is a first-class graph node with the following structure:

\begin{itemize}[nosep]
    \item \textbf{Label:} Human-readable name (\eg, ``cooking'',
    ``rust-learning''), generated by LLM or derived from prototype content.
    \item \textbf{Prototype node:} The most representative semantic node
    (highest corroboration count among members).
    \item \textbf{Centroid embedding:} Running average of member embeddings,
    updated incrementally as members are assigned.
    \item \textbf{Member count:} Number of assigned semantic nodes.
    \item \textbf{Parent ID (v0.2.0):} Optional reference to a parent
    category, enabling hierarchical nesting. Sub-categories inherit their
    parent's domain while representing finer-grained distinctions---forming
    an emergent is-a taxonomy (e.g., ``programming'' $\rightarrow$
    ``rust-async'', ``python-ml'').
    \item \textbf{Stability:} A score in $[0, 1]$ tracking how consistently
    the category survives transformation cycles. Stability increases
    asymptotically: $s' = s + 0.1 \cdot (1 - s)$.
\end{itemize}

Categories participate in the Hebbian graph via bidirectional
\texttt{MemberOf} links (weight~0.3), enabling spreading activation to
flow through category nodes. This produces \emph{cross-domain bridging}:
a query about ``cooking'' activates the cooking category, which in turn
activates all member semantic nodes---including those that might not match
the query embedding directly but are categorically related. When categories
form hierarchies, activation can propagate from a sub-category through its
parent to sibling sub-categories, bridging knowledge domains that share a
common ancestor but have no direct Hebbian links between their members.

\subsection{Incremental Assignment (Consolidation Phase)}

When \texttt{consolidate()} creates new semantic nodes, each node is
evaluated for category membership using a two-stage process:

\begin{enumerate}[nosep]
    \item \textbf{Embedding signal:} If the node has an embedding, compute
    cosine similarity to each category's centroid. The best match above a
    threshold of 0.6 is selected.
    \item \textbf{Graph signal:} If no embedding match is found, examine the
    node's graph neighborhood (source episodes' linked semantic nodes). If a
    majority ($>$50\%) share a category, assign to that category.
\end{enumerate}

On assignment, the category's centroid is updated as a running average:
$\mathbf{c}' = \frac{n \cdot \mathbf{c} + \mathbf{v}}{n + 1}$, where
$\mathbf{c}$ is the current centroid, $\mathbf{v}$ is the new member's
embedding, and $n$ is the member count. A \texttt{MemberOf} link is created
in the Hebbian graph.

Critically, consolidation never creates new categories. This prevents
category proliferation from sparse early data and ensures that category
creation requires the corroborating evidence of multiple clustered nodes.

\subsection{Category Discovery (Transformation Phase)}

During \texttt{transform()}, after deduplication and pruning, the system
discovers new categories from uncategorized semantic nodes:

\begin{enumerate}[nosep]
    \item Gather all uncategorized nodes with embeddings. If fewer than
    three exist, skip discovery (minimum cluster size).
    \item Compute pairwise cosine similarity.
    \item Apply single-linkage agglomerative clustering with a merge
    threshold of 0.7.
    \item For each cluster with $\geq$3 members:
    \begin{enumerate}[nosep]
        \item Select the prototype node (highest corroboration count).
        \item Compute centroid embedding (mean of member embeddings).
        \item Generate label: LLM summarization of top member contents if
        available, otherwise first three words of prototype content.
        \item Create the category, assign members, create \texttt{MemberOf}
        links.
    \end{enumerate}
\end{enumerate}

\subsection{Category Maintenance}

Transformation also maintains existing categories through four operations:

\begin{itemize}[nosep]
    \item \textbf{Stability tracking:} Non-empty categories that survive a
    transform cycle have their stability incremented. Stability converges
    asymptotically to 1.0, reflecting increasing confidence in the
    category's validity.
    \item \textbf{Merge:} Categories whose centroids exceed 0.85 cosine
    similarity are merged---the higher-stability category absorbs the
    lower-stability one, combining members and recomputing the centroid.
    \item \textbf{Dissolution:} Categories with stability below 0.1
    are dissolved---members are unassigned and returned to the uncategorized
    pool for potential re-clustering.
    \item \textbf{Garbage collection:} Empty categories (all members removed
    by forgetting or deduplication) are deleted.
\end{itemize}

This maintenance cycle ensures that the ontology remains a living reflection
of the current knowledge state rather than accumulating stale categories.

\subsection{Hierarchical Splitting (v0.2.0)}

As knowledge accumulates, categories may grow to encompass increasingly
diverse members. When a category reaches 8 or more members and its internal
coherence drops below~0.6 (measured as the average cosine similarity of
member embeddings to the centroid), \texttt{transform()} triggers automatic
splitting:

\begin{enumerate}[nosep]
    \item Compute pairwise cosine similarity among all member embeddings.
    \item Apply union-find clustering with the same merge threshold used for
    discovery (0.7 cosine similarity).
    \item For each resulting cluster with $\geq$2 members, create a
    sub-category with \texttt{parent\_id} pointing to the original category.
    \item Reassign members to their sub-categories, updating
    \texttt{MemberOf} links.
    \item The original category remains as the parent node in the hierarchy.
\end{enumerate}

This produces a tree that refines as knowledge grows. A ``programming''
category may split into ``rust-systems'' and ``python-data-science''
sub-categories, each with its own centroid and stability trajectory.
The hierarchy is unbounded in depth---a sub-category that later grows
large and incoherent will itself split, producing increasingly granular
taxonomic structure without manual intervention.

The splitting mechanism complements the existing merge operation: merging
collapses over-fragmented categories while splitting decomposes
over-generalized ones. Together, they maintain an ontology that tracks
the natural granularity of the agent's knowledge.
